\section{Atomic descent and finite-pool inexactness}
\label{sec:descent}

This section converts the distinctions of \Cref{sec:setting} into energy inequalities. The standard
atomic residual estimate is derived first, since its orthogonality hypothesis is weakened in the
finite-tolerance algorithm.

\begin{lemma}[Atomic residual bound]
\label{lem:atomic-residual}
Under \Cref{ass:atomic} and the exact orthogonality \cref{eq:orthogonality},
\begin{equation}
  \norm{\E'(G_m)}_{\D^*}\geq \frac{\delta_m}{B}.
\label{eq:atomic-residual}
\end{equation}
\end{lemma}

\begin{proof}
Convexity and \(G_m\in\Phi_m\) give
\[
 \delta_m\leq\ip{\E'(G_m)}{G_m-u^*}
 =-\ip{\E'(G_m)}{u^*}.
\]
Approximate \(u^*/B\in K_1(\D)\) by absolute convex combinations of atoms. The pairing of each
combination with \(\E'(G_m)\) is bounded by the dictionary dual seminorm, and passage to the limit gives
\cref{eq:atomic-residual}.
\end{proof}

\begin{definition}[\(p\)-growth density]
\label{def:p-growth}
Let \(p\geq2\), write \(z=(s,\xi)\in\R\times\R^d\), and omit the scalar component \(s\) for a
pure-gradient energy. A twice differentiable convex density \(\Psi(x,z)\) has \(p\)-growth on the
relevant range if, for almost every \(x\),
\begin{align*}
 c_0\abs z^p-c_1&\leq \Psi(x,z)\leq c_2(1+\abs z^p),\\
 \norm{D_z^2\Psi(x,z)}&\leq c_3(1+\abs z)^{p-2},
\end{align*}
with positive constants independent of \(x\) and \(z\) in that range. The first line controls bounded
energy sublevels, whereas the Hessian bound is the part used in the directional estimate below.
\end{definition}

\begin{assumption}[Directional quadratic upper bound]
\label{ass:directional}
There is \(L>0\) such that, on the trajectory sublevel set, for every \(g\in\D\) and
\(\abs\lambda\leq1\),
\begin{equation}
 \E(u+\lambda g)\leq\E(u)+\lambda\ip{\E'(u)}g+\frac L2\lambda^2.
\label{eq:directional-smoothness}
\end{equation}
\end{assumption}

The assumption is directional; it does not require a globally Lipschitz derivative on the full Banach
space. For the standard \(p\)-energies, bounded \(W^{1,p}\) sublevels and an \(X\)-normalized dictionary
already give the needed estimate by H\"older's inequality; an \(L^\infty\) atom bound is only a
convenient stronger condition for more general densities. Indeed, for
\(F(\lambda)=\E(u+\lambda g)\), the diffusion part of \(F''\) is controlled by
\[
 C_p\int_\Omega(\abs{\nabla u}+\abs\lambda\abs{\nabla g}+\eps)^{p-2}
 \abs{\nabla g}^2\dd x,
\]
and H\"older's inequality gives
\(F''(\lambda)\leq C_p(\norm{\nabla u}_{L^p}+\norm{\nabla g}_{L^p}+\eps|\Omega|^{1/p})^{p-2}
\norm{\nabla g}_{L^p}^2\) for \(\abs\lambda\leq1\). A reaction term is handled similarly. Removing
almost-null projected atoms remains important for the numerical parameterization and for uniform
source/entropy constants, but it is not necessary for this abstract \(L^p\) estimate once the atoms are
already normalized in the energy space.

\begin{theorem}[Generic energy baseline]
\label{thm:generic}
Suppose \Cref{ass:atomic,ass:directional} hold and \(t_m\geq t>0\). Then the exact continuous CGA
satisfies
\begin{equation}
  \delta_m\leq \frac{C(B,L,t,\delta_0)}{m+1}.
\label{eq:generic-rate}
\end{equation}
\end{theorem}

\begin{proof}
The selected descent score is at least \(a_m=t\delta_m/B\) by \Cref{lem:atomic-residual}. Test the new
active space with \(G_m+\lambda g_{m+1}\) and choose
\(\lambda_m=\min\{1,a_m/L\}\). If \(a_m\leq L\), then
\begin{equation}
 \delta_{m+1}\leq\delta_m-\frac{t^2}{2LB^2}\delta_m^2.
\label{eq:quadratic-recurrence}
\end{equation}
If \(a_m>L\), the unit step reduces the energy by at least \(a_m/2>L/2\); such steps occur only
finitely often because the energy is bounded below. After this finite prefix, taking reciprocals in
\cref{eq:quadratic-recurrence} and summing proves \cref{eq:generic-rate}. The truncation at one is
essential because \Cref{ass:directional} is stated only on unit dictionary segments.
\end{proof}

\subsection{Inexact selection and correction}

Write \(\widehat\delta_m=\E(\widehat G_m)-\E(u^*)\). In place of exact orthogonality, suppose that
\begin{equation}
  \abs{\ip{\E'(\widehat G_m)}{\widehat G_m}}\leq\eta_m^{\rm stat}.
  \label{eq:orthogonality-defect}
\end{equation}
The decomposition contains one multiplicative resolution factor and three additive error
contributions: \(t_m^{\rm pool}\tau_m\) records weak selection and continuous-dictionary coverage,
\(\xi_m\) records score evaluation, \(\eta_m^{\rm stat}\) is the radial active-space stationarity
defect of \cref{eq:orthogonality-defect} alone, and \(\varepsilon_m^{\rm opt}\) records the energy
accuracy of the correction. The full active-space dual Karush--Kuhn--Tucker (KKT) residual is
denoted separately by \(r_m^{\mathrm{KKT}}\) below. In the terminology of the greedy-with-errors
framework of \citet[Section~3]{DereventsovTemlyakov2022}, \(t_m^{\rm pool}\tau_m\) is the effective
weakness parameter and \(\varepsilon_m^{\rm opt}\) is an energy-evaluation/correction error. The
quantities \(\xi_m\) and \(\eta_m^{\rm stat}\) separate the finite-score and biorthogonality
contributions that the abstract framework combines into a single stepwise inexactness allowance.

\begin{theorem}[Finite-pool one-step descent]
\label{thm:inexact}
Assume \Cref{ass:atomic,ass:directional}, \cref{eq:score-error,eq:pool-coverage}, and
\cref{eq:orthogonality-defect} on the numerical trajectory. Define
\[
 a_m=\left[
 t_m^{\rm pool}\tau_m\frac{\widehat\delta_m-\eta_m^{\rm stat}}{B}
 -(1+t_m^{\rm pool})\xi_m
 \right]_+.
\]
Here \([x]_+=\max\{x,0\}\).
Then
\begin{equation}
 \widehat\delta_{m+1}\leq\widehat\delta_m
 -\frac12\min\{a_m,a_m^2/L\}+\varepsilon_m^{\rm opt}.
 \label{eq:inexact-descent}
\end{equation}
If, for a fixed absorption constant \(\vartheta\in(0,1)\),
\begin{equation}
 \eta_m^{\rm stat}+\frac{B(1+t_m^{\rm pool})}{t_m^{\rm pool}\tau_m}\xi_m
 \leq\vartheta\widehat\delta_m,
 \label{eq:relative-defect}
\end{equation}
the effective weakness \(\inf_{m\ge m_0} t_m^{\rm pool}\tau_m\) is positive, and there
exist a finite index \(m_0\) and \(0\leq\zeta<1\) such that
\[
  \varepsilon_m^{\rm opt}\leq\zeta\,H_L(a_m),\qquad m\ge m_0,
\]
then \(\widehat\delta_m=O((m-m_0+1)^{-1})\) for the tail. The finite prefix
\(m<m_0\) is not included in this rate statement.
\end{theorem}

\begin{proof}
Convexity and \cref{eq:orthogonality-defect} give
\[
 \widehat\delta_m
 \leq\ip{\E'(\widehat G_m)}{\widehat G_m-u^*}
 \leq\eta_m^{\rm stat}+B\norm{\E'(\widehat G_m)}_{\D^*}.
\]
Combining this estimate with \Cref{lem:true-score} bounds the true score of the selected atom below by
\(a_m\). The inexact correction and the directional quadratic upper bound imply, for
\(0\leq\lambda\leq1\),
\[
 \widehat\delta_{m+1}\leq\widehat\delta_m-\lambda a_m+\frac{L}{2}\lambda^2+\varepsilon_m^{\rm opt}.
\]
Set \(\lambda=\min\{1,a_m/L\}\) to obtain \cref{eq:inexact-descent}. Under
\cref{eq:relative-defect}, \(a_m\) is a fixed positive fraction of \(\widehat\delta_m\). For
\(m\ge m_0\), the displayed \(\zeta\)-condition absorbs the correction error into the
principal decrease, leaving the recurrence used in \Cref{thm:generic} with modified constants.
\end{proof}

\begin{remark}[Error saturation]
\Cref{eq:inexact-descent} shows how error saturation arises. A pool whose
\(\tau_m\) decreases, a score error \(\xi_m\) comparable with the best resolvable score, or a correction
error \(\varepsilon_m^{\rm opt}\) comparable with the predicted decrease can each stop the recursion. The equation
does not identify which defect dominates from a plotted error curve alone.
\end{remark}

The optimization error can be related to a KKT residual only after a norm and a strong-convexity
constant have been fixed.

\begin{lemma}[KKT residual to correction error]
\label{lem:kkt}
Let \(\Phi\) be finite dimensional and let \(\E|_\Phi\) be \(\mu_\Phi\)-strongly convex in a norm
\(\norm\cdot_\Phi\). If \(G_\Phi^{\sharp}\) is the exact minimizer and
\[
 r^{\mathrm{KKT}}=\norm{\E'(\widehat G)|_\Phi}_{\Phi^*},
\]
then
\begin{equation}
 0\leq\E(\widehat G)-\E(G_\Phi^{\sharp})
 \leq\frac{(r^{\mathrm{KKT}})^2}{2\mu_\Phi}.
\label{eq:kkt-energy}
\end{equation}
The same residual also satisfies
\begin{equation}
 \abs{\ip{\E'(\widehat G)}{\widehat G}}
 \leq r^{\mathrm{KKT}}\norm{\widehat G}_\Phi.
\label{eq:kkt-orth}
\end{equation}
\end{lemma}

\begin{proof}
Strong convexity at \(\widehat G\) gives
\[
 \E(G_\Phi^{\sharp})\geq\E(\widehat G)+
 \ip{\E'(\widehat G)}{G_\Phi^{\sharp}-\widehat G}
 +\frac{\mu_\Phi}{2}\norm{G_\Phi^{\sharp}-\widehat G}_\Phi^2.
\]
Apply the dual inequality and optimize the resulting quadratic in the distance to obtain
\cref{eq:kkt-energy}; \cref{eq:kkt-orth} is the dual inequality with
\(v=\widehat G\).
\end{proof}
